\chapter{Introduction}



\section{Viral Vectors in Emerging Therapies}

Clinical medicine progressively transitioning from broad, non-specific treatments such as chemotherapy or immunosuppressants to more refined, high-fidelity solutions. Those include small molecules targeting pathways of the immune system, antibody treatments, and more recently, gene therapy and virotherapy. Virotherapy emerges as a potentially promising approach for treating various cancers \cite{xiaoOncolyticVirusesAdvanced2026}, while gene therapy shows promise for hereditary diseases \cite{shchaslyvyiCurrentStateHuman2023}, and genetic vaccination is being explored for the treatment and prevention of infectious diseases \cite{zhouDNABasedVaccinesAdvances2026}. Gene therapy utilizes a vector to deliver genetic material to the target cell. Numerous vector systems are currently in use and are broadly categorized into viral and non-viral vectors. Non-viral vectors are further categorized into physical and chemical vectors. Viral vectors are particularly favored due to their evolutionary adaptation for efficient cell transduction as intracellular parasites \cite{kajanVirusHostCoevolution2020}. The most prevalent viral vectors employed in clinical trials are adenoviral vectors, followed by adeno-associated virus and lentiviral vectors \cite{bulchaViralVectorPlatforms2021}.

\section{Adenoviruses as Vectors}

\section{Producer Cell Lines}
Test commit
\section{CRISPR sequences and Cas proteins}
In brief, the CRISPR-Cas-System can be described as either of two things: originally, it is a component of the bacterial arsenal to defend against foreign genetic material. After its discovery, it has been studied and employed to accurately and precisely induce strand breaks in DNA, as a means of facilitating genetic engineering. its components are the growing family of Cas (CRISPR-associated) nucleases which do the actual strand-nicking and the CRISPR-Sequences, "CRISPR" meaning \textit{Clustered Regularly Interspaced Short palindromic Repeats}. Long before the role of this system was understood, the repetitive interspaced sequences were discovered in \textit{Escherichia coli} \cite{ishinoNucleotideSequenceIap1987}. The discovery of similar sequences in archaea  sparked new interest in the topic, as the presence of similar sequences in phylogenetically distant organisms implies a high level of evolutionary conservation and as such a high biological importance \cite{Si2005}. It was later described, that the highly varied spacer sequences represent genetic information of foreign (viral) origin \cite{pourcelCRISPRElementsYersinia2005} \cite{bolotinClusteredRegularlyInterspaced2005} and are accompanied by various associated (Cas) genes \cite{jansenIdentificationGenesThat2002} that execute their function either in maintaining and expanding the CRISPR loci, called class I, or serving as direct effector molecules in the cleavage of foreign DNA or RNA - class II \cite{shmakovSystematicPredictionGenes2018}. At present time, over 30 different Cas genes of both classes and a multitude of subtypes have been described \cite{makarovaEvolutionaryClassificationCRISPR2020}. other essential components of the system are the CRISPR-RNA (crRNA) and the trans-activating crRNA (tracrRNA). the crRNA is complimentary in sequence to the targeted foreign DNA or RNA (the protospacer), is derived from the spacer region of the CRISPR locus and has a length of $\simeq$20bp. It directs the effector complex to its target \cite{brounsSmallCRISPRRNAs2008}. The tracrRNA is important for the maturation of the crRNA and assembly of the final functional effector ribnucleoproteincomplex of Cas protein and associated RNAs \cite{deltchevaCRISPRRNAMaturation2011}. additionally, the protospacer must be neighbouring an adjacent motif (PAM) varying in sequence depending on the Cas type and class. The PAM-motifs are not integrated into the spacer sequences but are recognized by the Cas proteins. For example the Cas9 of \textit{s. pyogenes} requires a NGG PAM sequence directly downstream of the crRNA binding site and cleaves 3 base pairs upstream of the PAM. Since crRNA and tracrRNA are two separate molecules, it was then proposed to combine the two into a single chimeric guide RNA (sgRNA) to facilitate easier application of the system \textit{in vitro} \cite{jinekProgrammableDualRNAguidedDNA2012}. \\\\ To apply the system to edit eukaryotic DNA, the Cas9 sequence is codon optimized and paired with a nuclear localization signal \cite{maliRNAGuidedHumanGenome2013}. For our application, genetically modifying HEK293 cells, we chose Cas9 as it targets DNA and is the most popular and well studied Cas system to date. 

\newpage

\section{The hit and run hypothesis}

This hypothesis is in contrast to the general consensus that persistent presence of viral gene products is necessary for maintaining a transformative cellular phenotype. It has first been proposed by Skinner et al. \cite{skinnerTransformationPrimaryHamster1976} in the context of HSV-2 mediated transformation of primary hamster fibroblasts. In the realm of adenovirus, ...



\begin{figure}[ht]
\image{WT_ad5 genome.png}
\caption{Schematic representation of the genome of adenovirus type 5}
\end{figure}





